Modern AI systems exercise sovereign power over individuals---granting
credit, denying healthcare, and influencing elections---yet they lack
the structural checks and balances that govern traditional institutions.
While technical solutions for transparency and accountability exist,
they are fragmented: type systems enforce invariants, logs provide
persistence, and zero-knowledge proofs offer privacy. No unified
framework exists to verify that these components compose into a
legitimate ``institution.''

This paper introduces \textbf{Constitutional Code}, a theoretical
framework that formalizes Montesquieu's \textit{Separation of Powers}
for algorithmic systems. We demonstrate a strict formal correspondence
between political theory and verified systems architecture:
(1)~\textbf{Legislative Power} maps to a Constitutional Type System,
where regulations are encoded as linear resource types that constrain
behaviour at compile-time and cannot be overridden by the operator at
runtime; (2)~\textbf{Executive Power} maps to the Operator, who
executes decisions but is structurally prevented from suppressing
evidence or altering the record; and (3)~\textbf{Judicial Power} maps
to Independent Auditors, who verify state integrity via Zero-Knowledge
proofs without requiring operator cooperation and without compromising
individual privacy.

We prove three independence theorems and a completeness corollary
showing that the capability sets of the three branches are mutually
disjoint---the first time Montesquieu's separation has been enforced
not by social contract but by the mathematical properties of the
execution substrate. We show that traditional ``best-effort''
architectures fail each independence condition, and we derive
operational implications for regulators: the shift from post-hoc
audit to ex-ante constitutional type certification. This framework
provides regulators with a rigorous tool to evaluate whether AI
systems are merely compliant or structurally just.
